\chapter{参考文献与交叉引用}

\section{文献引用方法}
在常规文本处理软件（如 Word）中，通常从学术检索平台复制预设格式的文本。但在 \LaTeX{} 系统中，该流程被标准化为数据驱动的模式：

\begin{figure}[htbp]
    \centering
    \includegraphics[width=0.8\textwidth]{image4.png}
    \caption{学术检索平台中文献引用接口示意图}
    \label{fig:baidu-scholar}
  \end{figure}

  \begin{enumerate}
    \item \textbf{数据提取}：如图~\ref{fig:baidu-scholar} 所示，应忽略顶部的格式化文本，选择底部的 \texttt{BibTeX} 导出项，获取结构化代码。
    \item \textbf{数据库维护}：将获取的结构化数据增补至项目根目录下的 \texttt{references.bib} 文件。
    \item \textbf{逻辑引用}：调用 \texttt{\textbackslash cite\{ID\}} 指令实现正文与数据库的关联。例如引用 Transformer 相关研究：\,\cite{transformer}。% “\,”可选，用于调整排版
\end{enumerate}

\remark{
论文提供的bibtex格式里面可能会自动加入DOI号，可以手动删除\textbf{bibtex}中的DOI字段，以保持和老版参考文献格式相同
}
